Let $\overline{\mathcal{M}}_{g,n}$ be the Deligne--Mumford
moduli space of genus $g$ complex stable algebraic curves
(possibly with nodes), with $n>0$ distinct labeled marked
points. Consider the tautological line bundles
$\mathcal{L}_i\to\overline{\mathcal{M}}_{g,n}$, $i=1,\dots,n$, defined fiberwise by
$\mathcal{L}_i|_{C,x_1,\dots,x_n} \cong
T^*_{x_i}C$, where $C$ is a~genus~$g$ curve with marked
points $x_1,\dots,x_n$ (the definition makes sense since
the marked points are not allowed to coincide with the nodes).

In the late 1980-ies E.~Witten \cite{Witten} introduced a
theory of two-dimensional topological gravity, where the
classes $\psi_i=c_1(\mathcal{L}_i)$, $i=1,\dots,n,$ played
the role of observables, and the intersection numbers
\begin{gather*}
\langle\tau_{d_1}\cdots\tau_{d_n}\rangle_{g,n}
=\big\langle\psi_1^{d_1}\cdots\psi_n^{d_n}\big\rangle_{g,n}
=\int_{\overline{\mathcal{M}}_{g,n}}
\psi_1^{d_1}\cdots\psi_n^{d_n} ,
\end{gather*}
where $d_1+\dots +d_n=3g-3+n$, represented correlators of
the theory. Following a common convention, we will omit $g$
and $n$, or just $n$ when they are clear from the context.
